% Options for packages loaded elsewhere
\PassOptionsToPackage{unicode}{hyperref}
\PassOptionsToPackage{hyphens}{url}
\documentclass[
]{article}
\usepackage{xcolor}
\usepackage{amsmath,amssymb}
\setcounter{secnumdepth}{-\maxdimen} % remove section numbering
\usepackage{iftex}
\ifPDFTeX
  \usepackage[T1]{fontenc}
  \usepackage[utf8]{inputenc}
  \usepackage{textcomp} % provide euro and other symbols
\else % if luatex or xetex
  \usepackage{unicode-math} % this also loads fontspec
  \defaultfontfeatures{Scale=MatchLowercase}
  \defaultfontfeatures[\rmfamily]{Ligatures=TeX,Scale=1}
\fi
\usepackage{lmodern}
\ifPDFTeX\else
  % xetex/luatex font selection
\fi
% Use upquote if available, for straight quotes in verbatim environments
\IfFileExists{upquote.sty}{\usepackage{upquote}}{}
\IfFileExists{microtype.sty}{% use microtype if available
  \usepackage[]{microtype}
  \UseMicrotypeSet[protrusion]{basicmath} % disable protrusion for tt fonts
}{}
\makeatletter
\@ifundefined{KOMAClassName}{% if non-KOMA class
  \IfFileExists{parskip.sty}{%
    \usepackage{parskip}
  }{% else
    \setlength{\parindent}{0pt}
    \setlength{\parskip}{6pt plus 2pt minus 1pt}}
}{% if KOMA class
  \KOMAoptions{parskip=half}}
\makeatother
\usepackage{color}
\usepackage{fancyvrb}
\newcommand{\VerbBar}{|}
\newcommand{\VERB}{\Verb[commandchars=\\\{\}]}
\DefineVerbatimEnvironment{Highlighting}{Verbatim}{commandchars=\\\{\}}
% Add ',fontsize=\small' for more characters per line
\newenvironment{Shaded}{}{}
\newcommand{\AlertTok}[1]{\textcolor[rgb]{1.00,0.00,0.00}{\textbf{#1}}}
\newcommand{\AnnotationTok}[1]{\textcolor[rgb]{0.38,0.63,0.69}{\textbf{\textit{#1}}}}
\newcommand{\AttributeTok}[1]{\textcolor[rgb]{0.49,0.56,0.16}{#1}}
\newcommand{\BaseNTok}[1]{\textcolor[rgb]{0.25,0.63,0.44}{#1}}
\newcommand{\BuiltInTok}[1]{\textcolor[rgb]{0.00,0.50,0.00}{#1}}
\newcommand{\CharTok}[1]{\textcolor[rgb]{0.25,0.44,0.63}{#1}}
\newcommand{\CommentTok}[1]{\textcolor[rgb]{0.38,0.63,0.69}{\textit{#1}}}
\newcommand{\CommentVarTok}[1]{\textcolor[rgb]{0.38,0.63,0.69}{\textbf{\textit{#1}}}}
\newcommand{\ConstantTok}[1]{\textcolor[rgb]{0.53,0.00,0.00}{#1}}
\newcommand{\ControlFlowTok}[1]{\textcolor[rgb]{0.00,0.44,0.13}{\textbf{#1}}}
\newcommand{\DataTypeTok}[1]{\textcolor[rgb]{0.56,0.13,0.00}{#1}}
\newcommand{\DecValTok}[1]{\textcolor[rgb]{0.25,0.63,0.44}{#1}}
\newcommand{\DocumentationTok}[1]{\textcolor[rgb]{0.73,0.13,0.13}{\textit{#1}}}
\newcommand{\ErrorTok}[1]{\textcolor[rgb]{1.00,0.00,0.00}{\textbf{#1}}}
\newcommand{\ExtensionTok}[1]{#1}
\newcommand{\FloatTok}[1]{\textcolor[rgb]{0.25,0.63,0.44}{#1}}
\newcommand{\FunctionTok}[1]{\textcolor[rgb]{0.02,0.16,0.49}{#1}}
\newcommand{\ImportTok}[1]{\textcolor[rgb]{0.00,0.50,0.00}{\textbf{#1}}}
\newcommand{\InformationTok}[1]{\textcolor[rgb]{0.38,0.63,0.69}{\textbf{\textit{#1}}}}
\newcommand{\KeywordTok}[1]{\textcolor[rgb]{0.00,0.44,0.13}{\textbf{#1}}}
\newcommand{\NormalTok}[1]{#1}
\newcommand{\OperatorTok}[1]{\textcolor[rgb]{0.40,0.40,0.40}{#1}}
\newcommand{\OtherTok}[1]{\textcolor[rgb]{0.00,0.44,0.13}{#1}}
\newcommand{\PreprocessorTok}[1]{\textcolor[rgb]{0.74,0.48,0.00}{#1}}
\newcommand{\RegionMarkerTok}[1]{#1}
\newcommand{\SpecialCharTok}[1]{\textcolor[rgb]{0.25,0.44,0.63}{#1}}
\newcommand{\SpecialStringTok}[1]{\textcolor[rgb]{0.73,0.40,0.53}{#1}}
\newcommand{\StringTok}[1]{\textcolor[rgb]{0.25,0.44,0.63}{#1}}
\newcommand{\VariableTok}[1]{\textcolor[rgb]{0.10,0.09,0.49}{#1}}
\newcommand{\VerbatimStringTok}[1]{\textcolor[rgb]{0.25,0.44,0.63}{#1}}
\newcommand{\WarningTok}[1]{\textcolor[rgb]{0.38,0.63,0.69}{\textbf{\textit{#1}}}}
\setlength{\emergencystretch}{3em} % prevent overfull lines
\providecommand{\tightlist}{%
  \setlength{\itemsep}{0pt}\setlength{\parskip}{0pt}}
\usepackage{fancyhdr}
\usepackage{eso-pic}
\usepackage{tikz}
\usepackage{geometry}
\usepackage{xcolor}
\geometry{margin=1in}
\AddToShipoutPictureBG{%
  \begin{tikzpicture}[remember picture, overlay]
    \node[rotate=45, scale=3.5, text opacity=0.10, color=gray]
      at (current page.center)
      {\textsf{Gem v0.1.0 \textperiodcentered\ 2026-04-16}};
  \end{tikzpicture}%
}
\pagestyle{fancy}
\fancyhf{}
\fancyhead[L]{\small\textbf{Gem Language} v0.1.0}
\fancyhead[R]{\small 2026-04-16}
\fancyfoot[C]{\small\thepage}
\renewcommand{\headrulewidth}{0.4pt}
\renewcommand{\footrulewidth}{0.4pt}
\usepackage{bookmark}
\IfFileExists{xurl.sty}{\usepackage{xurl}}{} % add URL line breaks if available
\urlstyle{same}
\hypersetup{
  hidelinks,
  pdfcreator={LaTeX via pandoc}}

\title{Gem Language Reference}
\author{David B Hon}
\date{2026-04-16}

\begin{document}
\maketitle

\section{David B Hon Hard Science Gem Language Copywrite 2025,
2026}\label{david-b-hon-hard-science-gem-language-copywrite-2025-2026}

\textbf{Version:} 0.1.0 \textbar{} \textbf{Date:} 2026-04-16 \textbar{}
\textbf{Built with:} Gemini CLI \& Kiro CLI

Welcome to the \textbf{Gem Language} --- a modern, expressive STEM
language built by Gemini-CLI and Kiro-CLI using C++26 for AI, scientific
computing, quantitative finance, and cross-platform application
development.

\subsection{Key Features}\label{key-features}

\begin{itemize}
\tightlist
\item
  \textbf{Explicit Typing with Initial Values}: All variable
  declarations require an initial value (\texttt{int\ x\ =\ 0}).
\item
  \textbf{Global vs.~Local Scope}: Variables starting with \texttt{\_}
  (e.g., \texttt{\_config}) are global; all others are local.
\item
  \textbf{Universal Inheritance}: Every object inherits from
  \texttt{sys}, giving immediate access to \texttt{sys.print}, etc.
\item
  \textbf{First-Class AI Support}: Built-in \texttt{ai} object supports
  Gemini, Mistral, and Ollama.
\item
  \textbf{Polyglot Interop}: \texttt{use} keyword bridges Python, Julia,
  R, Fortran, C++, Go, Ruby, Rust, Node.js on-the-fly.
\item
  \textbf{Mobile \& Cross-Platform}: \texttt{mobl} object + browser PWA
  (Web Speech + Geolocation) works on Android, iPhone, macOS, Linux, and
  Windows 11.
\item
  \textbf{Domain-Specific Built-ins}: \texttt{fin}, \texttt{bsm},
  \texttt{bev}, \texttt{geo}, \texttt{data}, \texttt{chart},
  \texttt{mobl}, \texttt{trek}, \texttt{astro}, \texttt{nlp}, \texttt{www}, \texttt{cdn}, \texttt{thread} built into the runtime.
\item
  \textbf{Regular Expressions}: \texttt{rex} builtin provides full
  ECMAScript regex --- match, find, findall, groups, sub, gsub, split,
  count.
\item
  \textbf{Symbolic Math}: \texttt{math} builtin supports symbolic
  differentiation, integration, simplification, and LaTeX output via
  SymPy/Sage.
\item
  \textbf{Astrophysics}: \texttt{astro} builtin covers stellar physics,
  orbital mechanics, cosmology, solar physics, and exoplanets.
\end{itemize}

\begin{center}\rule{0.5\linewidth}{0.5pt}\end{center}

\subsection{Tutorial Roadmap}\label{tutorial-roadmap}

\subsubsection{1. The Foundations}\label{the-foundations}

\begin{itemize}
\tightlist
\item
  \textbf{\url{01_basics.g}}: Variables, types, and basic I/O using the
  \texttt{sys} object.
\item
  \textbf{\url{02_math_latex.g}}: Arithmetic operations and LaTeX-style
  mathematical expressions.
\item
  \textbf{\url{03_text_pdf.g}}: Text manipulation and PDF document
  generation.
\item
  \textbf{\url{04_images.g}}: Image processing with ImageMagick.
\item
  \textbf{\url{05_geo_gis.g}}: GIS, geolocation, GeoJSON, and
  interactive maps.
\end{itemize}

\subsubsection{2. Application
Development}\label{application-development}

\begin{itemize}
\tightlist
\item
  \textbf{\url{06_web_apps.g}}: Building web servers with
  \texttt{sys.app()} and route handlers.
\item
  \textbf{\url{07_ai_nlp.g}}: NLP and AI prompting with Gemini, Mistral,
  and Ollama.
\item
  \textbf{\url{08_functions.gm}}: Defining functions, objects, and
  inheritance.
\item
  \textbf{\url{09_module_demo.g}}: Structuring code with modules.
\item
  \textbf{\url{11_modules_paths.g}}: Managing module resolution and
  import paths.
\end{itemize}

\subsubsection{3. Advanced Integration \&
AI}\label{advanced-integration-ai}

\begin{itemize}
\tightlist
\item
  \textbf{\url{10_interop_cpp.g}}: High-performance C++ interoperability
  via Cling JIT.
\item
  \textbf{\url{12_lib_management.g}}: Dependency management and library
  linking.
\item
  \textbf{\url{13_dependency_details.g}}: Deep dive into Gem's package
  ecosystem.
\item
  \textbf{\url{14_ai_providers.g}}: Using Gemini, Mistral, and Ollama
  side-by-side.
\item
  \textbf{\url{15_mistral_native.g}}: Native Mistral C++ bridge via
  \texttt{prompt\_native}.
\end{itemize}

\subsubsection{4. Specialized Tooling}\label{specialized-tooling}

\begin{itemize}
\tightlist
\item
  \textbf{\url{16_algo_text_extras.g}}: Advanced string algorithms and
  text analysis.
\item
  \textbf{\url{17_repl_features.g}}: Mastering the interactive Gem REPL.
\item
  \textbf{\url{18_tcp_sockets.g}}: Low-level networking fundamentals.
\item
  \textbf{\url{19_tcp_server.g}} \& \textbf{\url{20_tcp_client.g}}:
  Client-server architectures.
\item
  \textbf{\url{21_polyglot.g}}: Bridging Python, Julia, and Gem.
\end{itemize}

\subsubsection{5. Science, Finance \&
DevOps}\label{science-finance-devops}

\begin{itemize}
\tightlist
\item
  \textbf{\url{22_builtins_itr_tcp.g}}: Iterators and network protocols.
\item
  \textbf{\url{23_data_science.g}}: DataFrames, statistics, and
  visualization.
\item
  \textbf{\url{24_devops.g}}: Automation scripts and infrastructure
  management.
\item
  \textbf{\url{25_ollama_mistral.g}}: Running local LLMs via Ollama.
\item
  \textbf{\url{26_devops_containers_vms.g}}: Docker and Vagrant
  environments.
\end{itemize}

\subsubsection{6. High-Performance \& Symbolic
Math}\label{high-performance-symbolic-math}

\begin{itemize}
\tightlist
\item
  \textbf{\url{27_rust_node.g}}: Polyglot support for Rust and Node.js.
\item
  \textbf{\url{28_interactive_charts.g}}: Dynamic 2D/3D/multi-axes
  Plotly visualization.
\item
  \textbf{\url{29_quantlib_finance.g}}: Quantitative finance with
  QuantLib.
\item
  \textbf{\url{30_bsm_pde_inheritance.g}}: Black-Scholes-Merton PDE
  solvers.
\item
  \textbf{\url{31_history_and_compilation.g}}: Gem's compilation model
  and history.
\item
  \textbf{\url{32_symbolic_math.g}}: Symbolic differentiation and
  integration.
\item
  \textbf{\url{33_langport_porting.g}}: Porting legacy codebases to Gem.
\item
  \textbf{\url{34_geo_tectonics.g}}: Advanced GIS, tectonic modeling,
  and geo-visualization.
\item
  \textbf{\url{35_text_advanced.g}}: Large-scale text processing and
  indexing.
\end{itemize}

\subsubsection{7. Mobile \& Cross-Platform}\label{mobile-cross-platform}

\begin{itemize}
\tightlist
\item
  \textbf{\url{36_mobl_travel_log.g}}: \texttt{mobl} object --- NLP
  dictation + GPS + GeoJSON travel log PWA.
\end{itemize}

\subsubsection{8. System Programming \&
Hardware}\label{system-programming-hardware}

\begin{itemize}
\tightlist
\item
  \textbf{\url{37_device_drivers.g}}: \texttt{drvr} builtin --- Linux,
  Windows 11, macOS, and Android device driver development.
\end{itemize}

\subsubsection{9. Text \& Pattern Matching}\label{text-pattern-matching}

\begin{itemize}
\tightlist
\item
  \textbf{\url{38_rex.g}}: \texttt{rex} builtin --- full ECMAScript
  regular expressions: match, find, findall, groups, sub, gsub, split,
  count.
\end{itemize}

\subsubsection{10. Travel Logs \& OSM Mapping}\label{travel-logs-osm-mapping}

\begin{itemize}
\tightlist
\item
  \textbf{\url{39_trek.g}}: \texttt{trek} builtin --- create, display,
  and edit travel logs on OpenStreetMap with GeoJSON/GPX support.
\end{itemize}

\begin{center}\rule{0.5\linewidth}{0.5pt}\end{center}

\subsection{Help \& Builtins}\label{help-builtins}

\begin{Shaded}
\begin{Highlighting}[]
\NormalTok{{-}{-}{-} Gem Language Help {-}{-}{-}}
\NormalTok{Gem is a high{-}performance STEM language implemented in C++26.}

\NormalTok{Usage Tips:}
\NormalTok{  {-} All variables must be initialized: \textquotesingle{}int x = 0\textquotesingle{}}
\NormalTok{  {-} Global variables start with underscore: \textquotesingle{}\_config = 1\textquotesingle{}}
\NormalTok{  {-} Use \textquotesingle{}use "file.py"\textquotesingle{} to translate and run foreign code.}
\NormalTok{  {-} Use \textquotesingle{}langport("*.py", "out.gm")\textquotesingle{} to port and save code.}
\NormalTok{  {-} Use \textquotesingle{}lib\textquotesingle{} to see all loaded modules.}
\NormalTok{  {-} Use \textquotesingle{}his\textquotesingle{} to view today\textquotesingle{}s session history.}

\NormalTok{CLI Options:}
\NormalTok{  gem \textless{}file.g\textgreater{}              {-} Run a script.}
\NormalTok{  gem {-}c \textless{}main.g\textgreater{} [mod...]  {-} Compile multiple files into a binary.}
\NormalTok{  gem {-}o \textless{}name\textgreater{}             {-} Specify output name for compiler.}
\NormalTok{  gem {-}h                    {-} Print history and exit.}
\NormalTok{  gem {-}t \textless{}file\textgreater{} [{-}o output] {-} AI{-}assisted translation to Gem.}

\NormalTok{Available Builtin Modules:}
\NormalTok{  sys, math, ai, text, rex, algo, bev, file, zip, nlp, img, www, cdn, geo,}
\NormalTok{  cpp, tcp, itr, thread, data, k3s, vm, go, ruby, node, rust,}
\NormalTok{  fin, bsm, chart, astro, mobl, trek, drvr}

\NormalTok{Keywords for Documentation:}
\NormalTok{  fun, obj, use, alias, his, lib, end, if, while,}
\NormalTok{  int, double, string, bool, exit}

\NormalTok{Mobile \& Cross{-}Platform:}
\NormalTok{  mobl phone("name")  phone.dictate(text)  phone.make\_feature(lat,lon,text)}
\NormalTok{  ./gem travel\_log.g  ->  http://localhost:8080  (Android/iPhone/macOS/Linux/Win11)}

\NormalTok{Detailed help: help "topic"  or  help(topic)}
\end{Highlighting}
\end{Shaded}

\begin{center}\rule{0.5\linewidth}{0.5pt}\end{center}

\subsection{Builtin Module Reference}\label{builtin-module-reference}

\subsubsection{sys --- System Root Module}\label{sys-system-root-module}

Every object inherits from \texttt{sys}. Methods:
\begin{itemize}
\tightlist
\item \texttt{sys.print(...)} --- print to stdout
\item \texttt{sys.args()} --- command-line arguments as a vector
\item \texttt{sys.async(func, [params], [timeout])} --- run function in background thread
\item \texttt{sys.sethandler(signals, func)} --- set custom signal handler
\item \texttt{sys.host()} --- OS information string
\item \texttt{sys.exec(cmd)} --- execute a shell command
\item \texttt{sys.doc([path])} --- read file comments or current script comments
\item \texttt{sys.help([topic])} --- interactive help
\item \texttt{sys.exit()} --- exit the REPL or script
\item \texttt{sys.langport(pattern, [output])} --- port foreign code (e.g.\ \texttt{*.py}) to Gem
\item \texttt{sys.redirect(url, [port])} --- HTTP redirect or background redirect server
\item \texttt{sys.app(port, [routes])} --- background web app server; routes map paths to strings, files, or handler functions
\end{itemize}

\subsubsection{math --- Mathematics \& Symbolic Math}\label{math-mathematics-symbolic-math}

Standard functions: \texttt{sin(x)}, \texttt{cos(x)}, \texttt{sqrt(x)}, constant \texttt{math.pi}.

Symbolic math (SymPy/Sage backend):
\begin{itemize}
\tightlist
\item \texttt{math.useSymPy()}, \texttt{math.useSage()} --- switch symbolic provider
\item \texttt{math.simplify(expr)} --- symbolic simplification
\item \texttt{math.diff(expr, [var])} --- symbolic differentiation
\item \texttt{math.integrate(expr, [var])} --- symbolic integration
\item \texttt{math.solve(expr, [var])} --- solve algebraic equation
\item \texttt{math.sym\_latex(expr)} --- LaTeX for a symbolic expression
\item \texttt{math.to\_latex(val)} --- convert value or matrix to LaTeX string
\item \texttt{math.write\_latex(path, content, [standalone])} --- write LaTeX to file
\item \texttt{math.read\_latex(path)} --- read content from LaTeX document environment
\item \texttt{math.parse\_latex(text)} --- extract math expressions from text (\texttt{\$...\$})
\item \texttt{math.compile\_latex(path)} --- run pdflatex on a file
\end{itemize}

\subsubsection{ai --- AI Integration}\label{ai-integration}

Supports Gemini (default), Mistral, and Ollama:
\begin{itemize}
\tightlist
\item \texttt{ai.prompt(text)} --- send prompt to current provider
\item \texttt{ai.prompt\_native(text)} --- high-speed native Mistral C++ bridge
\item \texttt{ai.useMistral(model)}, \texttt{ai.useOllama(model, [host])}, \texttt{ai.useGemini()} --- switch provider
\item \texttt{ai.setKey(key)}, \texttt{ai.setHost(host)}, \texttt{ai.setPath(path)} --- API configuration
\item Properties: \texttt{provider}, \texttt{model}, \texttt{host}
\end{itemize}

\subsubsection{text --- String \& Document Processing}\label{text-string-document-processing}

\begin{itemize}
\tightlist
\item \texttt{text.read(path)} --- read file into string
\item \texttt{text.sub(s, start, [len])} --- substring
\item \texttt{text.replace(s, old, new)} --- replace all occurrences
\item \texttt{text.write\_pdf(path, text)} --- create PDF; \texttt{write\_pdf\_a} for PDF/A archival
\item \texttt{text.read\_pdf(path)} --- extract text from PDF
\item \texttt{text.read\_markdown(path)}, \texttt{text.write\_markdown(path, text)} --- Markdown
\item \texttt{text.read\_yaml(path)}, \texttt{text.write\_yaml(path, obj)} --- YAML
\item \texttt{text.read\_html(path)}, \texttt{text.write\_html(path, text)} --- HTML
\item \texttt{text.read\_xml(path)}, \texttt{text.write\_xml(path, text)} --- XML
\item \texttt{text.read\_fits\_header(path)} --- read FITS file header as object
\item \texttt{text.read\_hdf\_header(path)} --- read HDF5 file header as object
\end{itemize}

\subsubsection{geo --- GIS \& Geolocation}\label{geo-gis-geolocation}

\begin{itemize}
\tightlist
\item \texttt{geo.lookup()} --- IP-based geolocation (\texttt{.lat}, \texttt{.lon}, \texttt{.city}, \texttt{.country})
\item \texttt{geo.distance(lat1, lon1, lat2, lon2)} --- Haversine distance in km
\item \texttt{geo.write\_geojson(path, features)} --- write GeoJSON FeatureCollection
\item \texttt{geo.history(plate)} --- geological history of a tectonic plate (AI-assisted)
\item \texttt{geo.plot2d(data, [layout])} --- 2D Plotly map (open-street-map / scattermapbox)
\item \texttt{geo.plot3d(data, [layout])} --- 3D globe visualization (orthographic scattergeo)
\end{itemize}

\subsubsection{astro --- Astrophysics \& Planetary Science}\label{astro-astrophysics}

Constants: \texttt{G}, \texttt{c}, \texttt{AU}, \texttt{pc}, \texttt{ly}, \texttt{Msun}, \texttt{Rsun}, \texttt{Lsun}, \texttt{Tsun}, \texttt{Mearth}, \texttt{Rearth}, \texttt{H0}, \texttt{sigma\_sb}

Unit conversions: \texttt{to\_ly(pc)}, \texttt{to\_pc(ly)}, \texttt{to\_au(m)}, \texttt{deg\_to\_rad(d)}, \texttt{rad\_to\_deg(r)}

Stellar: \texttt{luminosity(R,T)}, \texttt{stefan\_boltzmann(T,R)}, \texttt{wien(T)}, \texttt{abs\_magnitude(m,d\_pc)},
\texttt{distance\_modulus(d\_pc)}, \texttt{spectral\_class(T)}, \texttt{schwarzschild\_radius(M)}, \texttt{escape\_velocity(M,R)}

Orbital: \texttt{orbital\_period(a\_au,M\_msun)}, \texttt{orbital\_velocity(a\_m,M\_kg)}, \texttt{hill\_sphere(a,mp,ms)},
\texttt{roche\_limit(R,rho\_p,rho\_s)}, \texttt{synodic\_period(p1,p2)}, \texttt{planet(name)} $\to$ Mercury\ldots{}Neptune

Solar: \texttt{solar\_flux(dist\_au)}, \texttt{solar\_wind\_pressure(n\_cm3,v\_km\_s)}, \texttt{sunspot\_cycle(year)},
\texttt{parker\_spiral\_angle(v\_sw,dist\_au)}, \texttt{solar\_activity()}

Cosmology: \texttt{hubble\_distance(z)}, \texttt{redshift\_velocity(z)}, \texttt{lookback\_time(z)}, \texttt{critical\_density()}

Coordinates: \texttt{equatorial\_to\_galactic(ra,dec)}, \texttt{angular\_separation(ra1,dec1,ra2,dec2)}

Exoplanet: \texttt{transit\_depth(rp\_rearth,rs\_rsun)}, \texttt{habitable\_zone(L\_lsun)}, \texttt{equilibrium\_temp(L,d,A)}

Pulsar: \texttt{pulsar\_spindown(P\_s, Pdot)} $\to$ \texttt{\{age\_yr, Bfield\_G, edot\_W\}}

\subsubsection{trek --- Travel Logs \& OSM Mapping}\label{trek-travel-logs-osm}

\begin{itemize}
\tightlist
\item \texttt{trek.new(path)} --- create a new empty GeoJSON travel log file
\item \texttt{trek.add(path, lat, lon, [title], [note])} --- append a waypoint
\item \texttt{trek.edit(path, index, title, note)} --- update title/note of waypoint at index
\item \texttt{trek.remove(path, index)} --- remove waypoint at index
\item \texttt{trek.load(path)} --- return GeoJSON string
\item \texttt{trek.show(path, [port])} --- serve interactive OSM (Leaflet) map with live edit UI (default port 8090)
\item \texttt{trek.export\_gpx(geojson\_path, gpx\_path)} --- export waypoints to GPX format
\item \texttt{trek.stats(path)} --- returns object with \texttt{.waypoints} count and \texttt{.distance\_km} total
\end{itemize}

Map tiles: OpenStreetMap (Leaflet.js). Waypoints stored as GeoJSON FeatureCollection. Edits persist to file.

\subsubsection{mobl --- Mobile \& Cross-Platform PWA}\label{mobl-mobile}

\begin{Shaded}
\begin{Highlighting}[]
\NormalTok{mobl phone("device\_name")}
\NormalTok{phone.dictate(spoken\_text)         \# NLP parse -> JSON \{title, note, tags\}}
\NormalTok{phone.make\_feature(lat, lon, text) \# GPS + dictation -> GeoJSON Feature}
\end{Highlighting}
\end{Shaded}

Browser supplies GPS (Geolocation API) + speech (Web Speech API). Gem server handles NLP, GeoJSON assembly, and file persistence.

Routes: \texttt{/} $\to$ PWA HTML, \texttt{/log} $\to$ POST \texttt{\{lat,lon,text\}}, \texttt{/data} $\to$ GET GeoJSON FeatureCollection

\subsubsection{fin --- Financial Engineering}\label{fin-financial-engineering}

\begin{itemize}
\tightlist
\item \texttt{fin.ticker(symbol)} --- real-time market data via yfinance
\item \texttt{fin.high\_yield\_bonds()}, \texttt{fin.high\_yield\_etfs()}, \texttt{fin.high\_yield\_equities()} --- top 50 from TradingView
\item \texttt{fin.bs\_price(type, strike, spot, rate, vol, t)} --- European option pricing
\item \texttt{fin.greeks(type, strike, spot, rate, vol, t)} $\to$ NPV, Delta, Gamma, Theta, Vega
\item \texttt{fin.date(d,m,y)}, \texttt{fin.calendar(name)}, \texttt{fin.is\_holiday(cal,d,m,y)}, \texttt{fin.add\_days(...)}, \texttt{fin.diff\_days(...)}
\end{itemize}

\subsubsection{bsm --- Black-Scholes-Merton}\label{bsm-black-scholes-merton}

Inherits all \texttt{fin} methods, plus:
\begin{itemize}
\tightlist
\item \texttt{bsm.price\_american(symbol, type, duration)} --- PDE-enhanced American option pricing (\texttt{weekly}/\texttt{monthly}/\texttt{quarterly})
\end{itemize}

\subsubsection{chart --- Interactive Plotting (Plotly.js)}\label{chart-interactive-plotting}

\begin{itemize}
\tightlist
\item \texttt{chart.plot(traces, [layout])} --- generate interactive HTML chart
\item \texttt{chart.show(path)} --- open chart in browser
\item \texttt{chart.server(path, [port])} --- serve chart via background HTTP server
\end{itemize}

\subsubsection{data --- Data Science}\label{data-data-science}

\begin{itemize}
\tightlist
\item \texttt{data.read\_csv(path)} --- read CSV into vector of row objects
\item \texttt{data.mean(vector)} --- arithmetic mean
\item \texttt{data.std(vector)} --- sample standard deviation
\end{itemize}

\subsubsection{rex --- Regular Expressions (ECMAScript)}\label{rex-regular-expressions}

All functions accept optional flags string (\texttt{"i"} = case-insensitive):
\begin{itemize}
\tightlist
\item \texttt{rex.match(text, pattern, [flags])} $\to$ bool
\item \texttt{rex.find(text, pattern, [flags])} $\to$ first match string
\item \texttt{rex.findall(text, pattern, [flags])} $\to$ list of all matches
\item \texttt{rex.groups(text, pattern, [flags])} $\to$ capture groups from first match
\item \texttt{rex.sub(text, pattern, repl, [flags])} $\to$ replace first match
\item \texttt{rex.gsub(text, pattern, repl, [flags])} $\to$ replace all matches
\item \texttt{rex.split(text, pattern, [flags])} $\to$ split on pattern
\item \texttt{rex.count(text, pattern, [flags])} $\to$ count non-overlapping matches
\end{itemize}

\subsubsection{algo --- Algorithms \& Time}\label{algo-algorithms-time}

\begin{itemize}
\tightlist
\item \texttt{algo.add(...)} --- sum all numeric arguments
\item \texttt{algo.quicksort(v)}, \texttt{algo.sort(v, [start], [end])} --- sort numeric vectors
\item \texttt{algo.now()} --- current local time as string
\item \texttt{algo.date\_add(ts, days)}, \texttt{algo.date\_diff(t1, t2)} --- date arithmetic
\end{itemize}

\subsubsection{tcp --- TCP/IP Networking}\label{tcp-networking}

\begin{itemize}
\tightlist
\item \texttt{tcp.listen(port)}, \texttt{tcp.accept(fd)}, \texttt{tcp.connect(host, port)} --- socket management
\item \texttt{tcp.send(fd, msg)}, \texttt{tcp.recv(fd, [size])}, \texttt{tcp.close(fd)} --- data transfer
\item \texttt{tcp.nic()} --- network interface configurations
\item \texttt{tcp.routes()} --- system routing table
\end{itemize}

\subsubsection{thread --- Background Threads}\label{thread-background-threads}

Returned by \texttt{sys.async()}. Methods on the thread handle:
\begin{itemize}
\tightlist
\item \texttt{thread.wait()} --- block until complete, returns result (or timeout error string)
\item \texttt{thread.is\_finished()} --- bool, non-blocking status check
\end{itemize}

\subsubsection{www --- Web Framework}\label{www-web-framework}

\begin{itemize}
\tightlist
\item \texttt{www.wget(url, file)} --- download file via curl
\item \texttt{www.app(port, [routes])} --- background HTTP server; routes map paths to HTML strings
\item \texttt{www.redirect(target, [port])} --- start redirect server or return redirect header string
\item \texttt{www.map2d(geojson, output)} --- render 2D map via Mapnik (requires HAS\_MAPNIK)
\end{itemize}

\subsubsection{cdn --- Caching Proxy Server}\label{cdn-caching-proxy-server}

Inherits all \texttt{www} methods, plus:
\begin{itemize}
\tightlist
\item \texttt{cdn.start(port, routes\_map)} --- start caching reverse-proxy; routes map path prefixes to origin URLs
\item \texttt{cdn.stats()} --- returns object with \texttt{.hits}, \texttt{.misses}, \texttt{.bytes}, \texttt{.cached\_items}
\item \texttt{cdn.purge([path])} --- evict one path or \texttt{"*"} to clear entire cache
\item \texttt{cdn.config(type, content)} --- generate nginx/apache config
\end{itemize}

\subsubsection{k3s --- Docker \& Kubernetes}\label{k3s-docker-kubernetes}

\begin{itemize}
\tightlist
\item \texttt{k3s.docker\_run(image, cmd)}, \texttt{k3s.docker\_ps()}, \texttt{k3s.docker\_build(path, tag)}, \texttt{k3s.docker\_stop(id)}
\item \texttt{k3s.k3s\_apply(yaml)}, \texttt{k3s.k3s\_get(resource)}, \texttt{k3s.k3s\_pods()}, \texttt{k3s.k3s\_nodes()}, \texttt{k3s.k3s\_logs(pod)}
\end{itemize}

\subsubsection{vm --- Vagrant VMs}\label{vm-vagrant-vms}

\texttt{vm.init(box)}, \texttt{vm.up()}, \texttt{vm.ssh(cmd)}, \texttt{vm.status()}, \texttt{vm.halt()}, \texttt{vm.destroy()}

\subsubsection{drvr --- Device Driver Development}\label{drvr-device-driver-development}

\begin{itemize}
\tightlist
\item \texttt{drvr.linux(name)}, \texttt{drvr.win11(name)}, \texttt{drvr.macos(name)}, \texttt{drvr.android(name)} --- generate driver templates
\item \texttt{drvr.build(target)}, \texttt{drvr.deploy(target)} --- cross-platform build and deploy
\end{itemize}

\subsubsection{Other Builtins}\label{other-builtins}

\begin{itemize}
\tightlist
\item \textbf{nlp} --- Natural Language Processing
\item \textbf{bev} --- Bevington data reduction: \texttt{data(x,y)}, \texttt{fit\_line()}, \texttt{param(idx)}
\item \textbf{file} --- \texttt{file.write(path, content)}, \texttt{file.exists(path)}
\item \textbf{zip} --- \texttt{zip.compress(src, archive)}, \texttt{zip.decompress(archive, dest)}
\item \textbf{img} --- \texttt{img.resize(src, width, height, dest)}
\item \textbf{cpp} --- \texttt{cpp.repl()}, \texttt{cpp.exec(code)} --- C++26 JIT via Cling
\item \textbf{itr} --- \texttt{itr.range(n)}, \texttt{itr.while(cond\_fun, body\_fun)}
\item \textbf{go}, \textbf{ruby}, \textbf{node}, \textbf{rust} --- polyglot \texttt{run(input)} + language-specific helpers
\end{itemize}

\begin{center}\rule{0.5\linewidth}{0.5pt}\end{center}

\subsection{Mobile Travel Log App}\label{mobile-travel-log-app}

\begin{Shaded}
\begin{Highlighting}[]
\ExtensionTok{./gem}\NormalTok{ travel\_log.g          }\CommentTok{\# start server on :8080}
\CommentTok{\# Android/iPhone → http://\textless{}host{-}ip\textgreater{}:8080}
\CommentTok{\# macOS/Linux/Win11 → http://localhost:8080}
\end{Highlighting}
\end{Shaded}

Uses the \texttt{mobl} builtin + \texttt{travel\_log.html} (PWA). Voice dictation
→ AI NLP → GPS pin → live GeoJSON → Plotly map.

\begin{center}\rule{0.5\linewidth}{0.5pt}\end{center}

\subsection{Getting Started}\label{getting-started}

\subsubsection{Prerequisites}\label{prerequisites}

\begin{itemize}
\tightlist
\item
  \textbf{macOS/Linux}: \texttt{g++} (supporting C++2b/C++26),
  \texttt{make}, \texttt{python3}, \texttt{libreadline-dev}.
\item
  \textbf{Windows 11}: \href{https://www.msys2.org/}{MSYS2} with
  \texttt{mingw-w64-x86\_64-gcc}, \texttt{make}, and
  \texttt{mingw-w64-x86\_64-python}. Ensure
  \texttt{C:\textbackslash{}msys64\textbackslash{}mingw64\textbackslash{}bin}
  is in your PATH.
\end{itemize}

\subsubsection{Build and Run}\label{build-and-run}

\begin{Shaded}
\begin{Highlighting}[]
\FunctionTok{make}\NormalTok{ all}
\ExtensionTok{./gem}\NormalTok{ tutorial/01\_basics.g}
\end{Highlighting}
\end{Shaded}

On Windows, the binaries will be produced as \texttt{gem.exe} and
\texttt{gem\_test.exe}.

\begin{center}\rule{0.5\linewidth}{0.5pt}\end{center}

\subsection{Windows Support}\label{windows-support}

Gem is now fully compatible with Windows 11 via MSYS2/MinGW-w64. Key
improvements include: - \textbf{Robust Dependency Detection}: Automatic
detection of Python headers and libraries using \texttt{sysconfig}. -
\textbf{Modern API Integration}: Support for \texttt{cpp-httplib} and
other libraries through \texttt{\_WIN32\_WINNT=0x0A00}. -
\textbf{Cross-Platform Signals}: Graceful handling of platform-specific
signals (\texttt{SIGINT}, \texttt{SIGTERM}). - \textbf{Path
Compatibility}: Seamless handling of Windows-style paths and
\texttt{.exe} extensions in compilation mode.

Or enter the REPL:

\begin{Shaded}
\begin{Highlighting}[]
\ExtensionTok{./gem}
\end{Highlighting}
\end{Shaded}

\begin{center}\rule{0.5\linewidth}{0.5pt}\end{center}

\emph{Developed with Gemini CLI and Kiro.}

\end{document}
